\section{System Architecture}
\label{sec:architecture}

Omni-Shield is a local-first document redaction system
designed around three principles:
(i)~no document or extracted PII ever leaves the
device;
(ii)~every redaction decision is explainable by the
agent that made it;
and (iii)~every completed redaction produces a
cryptographically verifiable audit record.
This section describes the overall system structure,
the nine-agent pipeline for image redaction, and the
multi-modal routing logic that extends the system to
five additional document formats.

% ─────────────────────────────────────────────────────────
\subsection{System Overview}
% ─────────────────────────────────────────────────────────

Figure~\ref{fig:architecture} shows the top-level
system structure.
The user interacts with a React/Electron frontend that
accepts document uploads, displays detected PII regions
on an interactive canvas editor, and exposes controls
for selective redaction (10 PII categories, custom
terms, regex patterns, and preset profiles).
A FastAPI backend running on \texttt{localhost:8000}
orchestrates all AI inference, file I/O, and blockchain
interaction.
A local Ganache Ethereum node on port 7545 maintains
the on-chain audit ledger.

All AI models run locally under 4-bit NF4
quantisation~\cite{dettmers2023qlora}:
Qwen2-VL-2B-Instruct (vision-language model, 1.40\,GB
VRAM) for document layout understanding,
Qwen2.5-3B-Instruct (text model, 1.91\,GB VRAM) for
contextual PII classification, EasyOCR for text
extraction, and YOLOv8n for face detection.
Total peak VRAM usage during image inference is
3.31\,GB, leaving headroom for OS and driver overhead
within the 8\,GB budget of the target hardware.
Models are loaded lazily on first request and persist
in VRAM for subsequent requests, reducing warm
inference latency to 10.42\,s mean for the full
image pipeline.

% ─────────────────────────────────────────────────────────
\subsection{AgentState: Shared Data Structure}
% ─────────────────────────────────────────────────────────

All agents communicate through a single shared
\texttt{AgentState} object that is passed sequentially
through the pipeline.
The state encapsulates the document path, extracted
text tokens with bounding boxes, VLM layout output,
detected PII items with pixel coordinates, and timing
metadata for each agent.
No agent communicates directly with any other agent;
all inter-agent communication is mediated by reading
and writing fields on the shared state.
This design makes the pipeline debuggable (every
intermediate state is inspectable), ablatable (any
agent can be disabled by skipping its state write),
and testable in isolation.

\begin{table}[h]\centering
\small
\caption{AgentState fields written by each agent}
\label{tab:agentstate}
\begin{tabular}{ll}
\toprule
Agent & Fields written \\
\midrule
RouterAgent      & \texttt{format}, \texttt{doc\_type}, \texttt{image\_path} \\
DeskewAgent      & \texttt{skew\_angle}, \texttt{image\_path} (corrected) \\
OCRAgent         & \texttt{ocr\_words} (text, bbox, confidence) \\
LayoutAgent      & \texttt{layout\_items} (label, value, value\_box), \texttt{doc\_type} \\
TextPIIAgent     & \texttt{text\_pii\_items} (entity, value) \\
ContextAgent     & \texttt{confirmed\_pii} (filtered, enriched) \\
BBRefinerAgent   & \texttt{confirmed\_pii} (pixel bboxes attached) \\
CriticAgent      & \texttt{confirmed\_pii} (false positives removed) \\
RedactionAgent   & \texttt{redacted\_path}, \texttt{coverage\_pct}, \texttt{face\_boxes} \\
\bottomrule
\end{tabular}
\end{table}

% ─────────────────────────────────────────────────────────
\subsection{The Nine-Agent Image Pipeline}
% ─────────────────────────────────────────────────────────

Image redaction (JPEG, PNG, BMP, TIFF) is the most
complex pathway, requiring all nine agents.
Figure~\ref{fig:pipeline} illustrates the data flow.

\paragraph{RouterAgent (0.04\,s mean).}
Inspects the uploaded file, determines its MIME type
and initial document type (\texttt{unknown},
\texttt{id\_card}, \texttt{passport}, etc.),
and routes to the appropriate pipeline variant.
For images, the full nine-agent pipeline is invoked.
For other formats, shortened pipelines are used
(Section~\ref{sec:multimodal}).

\paragraph{DeskewAgent.}
Estimates document skew using Hough line detection
and applies affine correction for angles exceeding
$\pm$0.5°.
Skew correction is applied before OCR to improve text
extraction accuracy on scanned documents.
The corrected image replaces the original in
\texttt{AgentState} for all downstream agents.

\paragraph{OCRAgent (2.73\,s mean).}
Runs EasyOCR with GPU acceleration to extract
all text tokens from the image.
Each token is stored with its bounding box in pixel
coordinates and confidence score.
OCR output serves two purposes: (i)~it provides the
text corpus for the TextPIIAgent and ContextAgent,
and (ii)~it enables the \emph{TEXT-FIRST} bounding
box refinement strategy in the BBRefinerAgent
(Section~\ref{sec:bbrefiner}).

\paragraph{LayoutAgent (4.79\,s mean).}
The LayoutAgent is the system's primary novel
component.
It submits the document image to the fine-tuned
Qwen2-VL-2B VLM with a structured prompt requesting
all PII label-value pairs with approximate bounding
boxes.
The VLM responds in a constrained format:

\vspace{4pt}
\noindent\texttt{LABEL: Name | VALUE: Angela Greene |
VALUE\_BOX: (200,180),(300,240)}
\vspace{4pt}

The LayoutAgent parses this output, applies label
normalisation (mapping variants such as
\texttt{SURNAME} and \texttt{Full Name} to the
canonical \texttt{Name} category), and additionally
uses the VLM's response to refine the document type
from \texttt{unknown} to a specific category
(e.g., \texttt{id\_card}).
The LayoutAgent accounts for 46\% of total pipeline
latency, reflecting the cost of VLM inference at
4-bit quantisation.

\paragraph{TextPIIAgent (0.94\,s mean).}
Runs the Qwen2.5-3B text model over the OCR-extracted
text to identify any PII entities that the VLM may
have missed --- particularly emails, phone numbers,
and free-form addresses that do not appear as
structured label-value pairs on the card.
This agent is most useful for document types with
dense running text (Word documents, plain text files)
and contributes marginally to identity card redaction
where text is sparsely structured.

\paragraph{ContextAgent (0.07\,s mean).}
Applies heuristic and LLM-based filtering to the
combined output of LayoutAgent and TextPIIAgent.
Heuristics encode domain knowledge: organisation
names, legal boilerplate, and card-issuer references
are removed even if the VLM labelled them as PII.
A lightweight LLM pass then reviews borderline cases,
using the user's selected PII filter (which of the
ten categories are enabled) to drop disabled categories
before downstream processing.

\paragraph{BBRefinerAgent (< 1\,ms mean).}
\label{sec:bbrefiner}
Refines the pixel-level bounding boxes for each
confirmed PII item using the \emph{TEXT-FIRST}
strategy: for each confirmed PII value, the agent
searches the OCR word list for a token whose
string matches the PII value and substitutes the
OCR-derived pixel box (which is precise to the
rendered glyph boundary) for the VLM-estimated box
(which is accurate to the nearest grid cell).
This substitution improves bounding box IoU for
text-valued PII categories without any additional
model inference.
Signature regions are handled separately via YOLOv8
contour detection, which produces full-width boxes
that encompass the drawn signature area.

\paragraph{CriticAgent (1.85\,s mean).}
Submits the confirmed PII list to the Qwen2.5-3B
text model for adversarial review.
The critic is instructed to identify items that are
not genuinely PII in context --- for example, a
nationality field on an ID card that happens to match
a person's name.
Items labelled as safe by the critic are dropped
\emph{unless} they carry a strong PII category label
(Name, Date of Birth, ID Number, Signature, Face),
in which case a \textit{PII label override} retains
them regardless of the critic's verdict.
This override prevents the critic from discarding
clearly sensitive information while still filtering
out false positives in ambiguous categories.

\paragraph{RedactionAgent (0.02\,s mean).}
Applies black rectangular redactions to the image
at the confirmed pixel bounding boxes using Pillow.
Face boxes detected by YOLOv8n are merged with the
VLM-detected PII list before redaction.
The agent records the total coverage percentage
(redacted pixels / total image area) and saves the
output image alongside a \texttt{.audit.json} sidecar
containing the full redaction record.

% ─────────────────────────────────────────────────────────
\subsection{Multi-Modal Routing}
\label{sec:multimodal}
% ─────────────────────────────────────────────────────────

Five additional document formats are handled by
shorter pipelines that reuse the hybrid NER component
(a combination of TextPIIAgent heuristics and
Qwen2.5-3B inference) but bypass the VLM.
Table~\ref{tab:latency} in Section~\ref{sec:evaluation}
reports latency for each format; we describe the
pipeline logic here.

\paragraph{PDF.}
Text is extracted natively via PyMuPDF without OCR.
The hybrid NER pipeline identifies PII entities,
and PyMuPDF's redaction API applies black rectangles
at the character-level bounding boxes of matched
text.
Mean latency is 40\,ms --- 225$\times$ faster than
image redaction --- because the VLM is not invoked.
PDF redaction is content-stream based and survives
copy-paste; the redacted text is genuinely absent
from the document, not merely obscured visually.

\paragraph{Audio.}
Whisper tiny~\cite{radford2022whisper} transcribes
the audio to text.
The hybrid NER pipeline identifies PII entity spans
in the transcript.
Span timestamps are mapped back to the audio timeline,
and pydub replaces those segments with silence.
Mean latency is 1.22\,s for a 7-second clip.

\paragraph{Plain text.}
The hybrid NER pipeline runs directly on the text
content.
Matched PII values are replaced with
\texttt{[REDACTED]} tags in the output.
Mean latency is 10\,ms.

\paragraph{Word documents (.docx).}
python-docx iterates all paragraphs and table cells,
applying string replacement for matched PII values.
Output is rendered as PDF with matched values
highlighted in red before redaction for user review.

\paragraph{Presentations (.pptx).}
python-pptx iterates all slide text frames.
Redaction follows the same string-replacement logic
as Word documents.

For all non-image formats, blockchain anchoring and
ZK proof generation (Section~\ref{sec:crypto})
proceed identically to the image pipeline.
The audit JSON sidecar records the document SHA-256,
anchor hash, redacted item list, coverage percentage,
and PII filter configuration used.

% ─────────────────────────────────────────────────────────
\subsection{Selective Redaction and User Controls}
% ─────────────────────────────────────────────────────────

Omni-Shield exposes fine-grained redaction controls
through the frontend.
Users select which of ten PII categories to redact
(names, faces, signatures, phones, emails, date of
birth, ID numbers, addresses, organisation names,
general dates), add arbitrary custom terms or regular
expressions, and choose from four preset profiles
(Full Scan, Legal, Medical, Minimal).

Custom terms and regex patterns are matched against
the OCR/NER output using exact string search and
Python \texttt{re} module respectively,
and are visually distinguished from AI-detected
redactions in the canvas editor (dark red borders
for custom, black for AI).
Dry-run mode executes the full pipeline and returns
the list of would-be redactions without writing any
output files or anchoring to the blockchain,
supporting preview workflows.
